\documentclass[conference]{IEEEtran}
\usepackage[utf8]{inputenc}
\usepackage{amsmath,amssymb}
\usepackage{booktabs}
\usepackage{graphicx}
\usepackage{xcolor}
\usepackage[colorlinks=true,linkcolor=blue,citecolor=blue,urlcolor=blue]{hyperref}

% ============================================================================
% IROS-WS DRIVING PAPER (rewritten 2026-07: metric pivot)
% Central metric changed from generation coherence to (i) graded NCAP closed-loop
% safety and (ii) free-drive route competence. Granularity demoted from "the fix"
% to a coherence-level finding. All numbers are real measurements
% (benchmark12 / freedrive / positive_hunt sweeps; 10 seeds; bootstrap CIs).
% ============================================================================

\title{\textbf{A Metric Hierarchy for Certifying Quantized Driving VLA Policies:\\
Open-Loop, Coherence, and Route Competence All Mislead;\\
Graded Closed-Loop Safety Decides}}
\author{
\IEEEauthorblockN{Justin Williams$^{1}$, Kishor Datta Gupta$^{2}$, Roy George$^{3}$, and Mrinmoy Sarkar$^{4}$
}\\
\IEEEauthorblockA{Department of Cyber-Physical Systems, Clark Atlanta University, Atlanta, GA, USA\\
$^{1}$\textit{justin.williams1@students.cau.edu},
$^{2}$\textit{kgupta@cau.edu},
$^{3}$\textit{rgeorge@cau.edu}}
\IEEEauthorblockA{Siemens Corporation, Princeton, NJ, USA\\
$^{4}$\textit{mrinmoy.sarkar@siemens.com}}
}

\begin{document}
\maketitle

\begin{abstract}
Vision--language--action (VLA) driving policies must be quantized before they can run on
in-vehicle hardware, and quantization is almost always validated on open-loop accuracy. We
show this validation is not merely optimistic but structurally blind---and that the natural
intermediate proxies are blind too. Evaluating post-training weight quantization of
OpenDriveVLA-0.5B with four metrics of increasing fidelity---open-loop L2, closed-loop
generation coherence, closed-loop route competence on an adversary-free benchmark we
construct, and graded closed-loop safety (impact-speed-graded NCAP score)---we find a strict
\emph{hierarchy of blindness}: each metric certifies quantized policies that the next
rejects. Open-loop L2 certifies everything (a $19\times$ ego-ablation shows it is
ego-extrapolation). Coherence certifies group-wise INT4, which restores well-formed
trajectories but not safety. Route competence certifies AWQ-INT4, which drives routes at
near-FP16 level ($87\%$ vs $94\%$) yet fails safety-critical scenarios exactly like naive
INT4. On the graded safety metric, \emph{every} 4-bit recipe we test---naive, group-wise,
AWQ, and DCT-rotated W4A4---collapses to $2.6$--$2.9$ of FP16's $4.4$ stars, with failures
concentrated on long-horizon stationary-obstacle scenarios; per-channel INT8 is statistically
indistinguishable from FP16 on both closed-loop benchmarks (paired $\Delta=-0.00$,
$95\%$ CI $[-0.17,+0.20]$) and is the recipe we certify. Single-scene closed-loop pilots
certified three of the failing recipes with perfect scores before full-scale evaluation
rejected them. All accuracy conclusions reproduce on a Jetson AGX Orin. We release the
harness, the free-drive benchmark, and a practitioner protocol: certify with the full
hierarchy, ship INT8, and treat sub-8-bit driving policies as uncertified.
\end{abstract}
\begin{IEEEkeywords}
vision-language-action models, autonomous driving, post-training quantization, closed-loop
evaluation, safety certification, NeuroNCAP, edge deployment, Jetson Orin
\end{IEEEkeywords}

\section{Introduction}
On-board inference for end-to-end driving stacks motivates compressing the language backbones
of driving VLAs, and the prevailing evidence---near-lossless 4-bit weight quantization on
open-loop nuScenes planning---suggests compression is essentially free. This paper shows that
the prevailing evidence is wrong in a specific, layered way: it is not only the open-loop
metric that is blind, but also the successively stronger closed-loop proxies one would
naturally reach for, and the recipes the LLM-quantization literature trusts (group-wise
scales, activation-aware scaling, orthogonal rotation) each pass some level of the hierarchy
while failing the one that matters.

Contributions:
\begin{itemize}
\item A \textbf{four-level metric hierarchy} for certifying quantized driving policies---
open-loop L2, generation coherence, route competence, graded closed-loop safety---with an
empirical demonstration that each level certifies configurations the next level rejects
(Table~\ref{tab:hierarchy}).
\item A controlled demonstration of open-loop blindness ($19\times$ ego-ablation), and
evidence that \emph{closed-loop} proxies short of graded safety are also insufficient:
coherence certifies group-wise INT4 and route competence certifies AWQ-INT4, both of which
fail graded safety at naive-INT4 level.
\item A \textbf{free-drive route-competence benchmark} on NeuroNCAP scenes (adversary-free,
cross-track-gated, anti-freeze-gated, with a collision-artifact taxonomy), which removes the
stay-in-place confound that contaminates adversarial scores of degenerate policies.
\item A full-scale verdict with confidence intervals: \textbf{per-channel INT8 is certified}
(indistinguishable from FP16 on both closed-loop benchmarks, $2\times$ compression), while
\textbf{no 4-bit weight recipe survives}; single-scene pilots certified three failing
recipes---a caution for closed-loop evaluation practice itself.
\item End-to-end \emph{on-Orin} systems and accuracy measurements showing the conclusions
transfer to deployment hardware.
\end{itemize}

\section{Related Work and Positioning}
\textbf{Driving planners and VLAs.} UniAD~\cite{uniad} introduced a planning-oriented,
query-based perception stack; VAD~\cite{vad} vectorized the scene for efficiency.
OpenDriveVLA~\cite{opendrivevla} reuses UniAD track+map perception as a vision tower for a
Qwen2.5 language/action head that emits trajectories as text, reaching state-of-the-art
open-loop L2 (Table~\ref{tab:planners}).
\textbf{Open-loop critique.} ``Is Ego Status All You Need?''~\cite{egostatus} showed nuScenes
open-loop planning is dominated by ego state; we quantify the effect ($19\times$) as it bears
on \emph{quantization} evaluation, and extend the critique to closed-loop proxies.
\textbf{Closed-loop neural simulation.} NeuroNCAP~\cite{neuroncap} couples a
NeuRAD~\cite{neurad} renderer to a safety-test protocol with an impact-speed-graded score;
we adopt that graded score as the top of our hierarchy and additionally construct an
adversary-free competence benchmark from the same scene reconstructions.
\textbf{LLM weight quantization.} RTN, GPTQ~\cite{gptq}, and AWQ~\cite{awq} rely on
group-wise scales to bound weight-outlier damage~\cite{llmint8,smoothquant}; rotation methods
(QuaRot~\cite{quarot}) suppress activation outliers. All are validated on language benchmarks
or open-loop accuracy. \textbf{Our positioning.} To our knowledge this is the first study to
evaluate driving-VLA quantization across a graded closed-loop safety metric, a closed-loop
competence metric, and on-device---and the first to show that the standard recipes pass the
cheaper metrics while failing graded safety.

\begin{table}[t]\centering
\caption{Open-loop nuScenes planning context (ST-P3 L2, m). Our contribution is not a new
planner but the observation that this metric cannot separate FP16 from quantized policies.}
\label{tab:planners}
\begin{tabular}{lcccc}
\toprule
Method & 1s & 2s & 3s & Avg. \\
\midrule
UniAD~\cite{uniad}            & 0.44 & 0.67 & 0.96 & 0.69 \\
VAD~\cite{vad}                & 0.17 & 0.34 & 0.60 & 0.37 \\
OpenDriveVLA-7B~\cite{opendrivevla}  & 0.15 & 0.31 & 0.55 & 0.33 \\
\textbf{OpenDriveVLA-0.5B} (study model) & 0.15 & 0.32 & 0.57 & 0.35 \\
\bottomrule
\end{tabular}
\end{table}

\section{Method}
\subsection{Model and quantization scope}
OpenDriveVLA-0.5B pairs Qwen2.5-0.5B (hidden $896$, $24$ layers) with UniAD stage-1 track+map
perception. We apply post-training quantization to the \textbf{168 Linear layers of the Qwen
backbone ($357.8$\,M parameters)}, leaving perception, projectors, embeddings, and the output
head in floating point, so observed effects are attributable to backbone precision. Weights
are dequantized to the compute dtype for inference (\emph{fake-quant}): the arithmetic path
is unchanged, so accuracy effects are real while packed memory gains are reported separately
(Sec.~\ref{sec:orin}).

\subsection{Quantized configurations}
\label{sec:configs}
Six configurations span the standard PTQ design space: \textbf{FP16} (reference);
\textbf{W8} ($b{=}8$, per-channel symmetric RTN); \textbf{naive W4} ($b{=}4$, per-channel);
\textbf{group W4} ($b{=}4$, group size $g{=}128$); \textbf{AWQ W4}~\cite{awq} ($b{=}4$,
$g{=}128$, activation-aware scales calibrated on a rendered-domain closed-loop rollout); and
\textbf{W4A4+DCT} (4-bit weights and activations under an orthonormal DCT rotation
$y=(WQ^\top)(Qx)$, QuaRot-style~\cite{quarot}). Group-wise scales cost $16/g=0.125$
bits/weight at $g{=}128$ ($\approx1\%$ footprint).

\paragraph{Why granularity should help---and what the argument actually proves.}
For a block $\mathcal{B}$ sharing scale
$s_\mathcal{B}=\max_{j\in\mathcal{B}}|w_j|/(2^{b-1}-1)$, any weight with
$|w_i|<\tfrac12 s_\mathcal{B}$ rounds to zero; at INT4 a single outlier
$w^\star$ erases every weight with $w^\star/|w_i|>14$ in its block. Per-channel blocks span
the whole row, so one outlier erases row-wide; group-wise confines the damage to $g$ weights.
This argument bounds \emph{weight-reconstruction} error, and it correctly predicts the
\emph{coherence} results (Sec.~\ref{sec:coherence}). It makes no prediction about closed-loop
\emph{safety}---and empirically (Sec.~\ref{sec:safety}) finer granularity does not restore
it. We state the mechanism at its true strength: necessary for well-formed generation, not
sufficient for safe driving.

\subsection{The four-level evaluation hierarchy}
\label{sec:hierarchy-def}
\textbf{(1) Open-loop L2} (standard): ST-P3 L2 on the $162$ nuScenes-mini$\,\cap\,$val
samples. The \emph{ego-ablation} zeroes ego-state and $2$\,s history in the prompt, leaving
perception untouched, to test how much of L2 is ego-extrapolation.
\textbf{(2) Generation coherence} (closed-loop): fraction of frames with a well-formed
(non-degenerate) trajectory during NeuroNCAP rollouts. A degenerate/malformed output triggers
a stay-in-place fallback.
\textbf{(3) Route competence} (closed-loop, ours): on \emph{adversary-free} rollouts of the
same NeuRAD scene reconstructions (original recorded traffic, no injected threat), a rollout
succeeds iff (i) maximum \emph{cross-track} distance to the logged route stays within
$\tau{=}4$\,m (the FP16 on-road envelope; cross-track rather than time-matched error, so a
correct path driven at a different speed is not penalized), (ii) driven path length exceeds
$5$\,m (\emph{anti-freeze gate}: a frozen policy fails rather than scoring ``safe''), and
(iii) no genuine collision. Collisions are classified by ego-frame geometry at impact:
\emph{spawn} (terminated at frame $\le1$; ego spawned in contact) and \emph{rear} (struck
from behind by logged traffic---the ego lawfully drives slower than the log) are simulation
artifacts and excluded; frontal/lateral strikes count.
\textbf{(4) Graded closed-loop safety}: the NeuroNCAP score---$5$ stars if no collision, else
$4\!\cdot\!\max(0,1-v_\text{impact}/v_\text{ref})$---graded by impact speed, so it
discriminates between full-speed and mitigated collisions even where the binary collision
rate saturates.

\subsection{Closed-loop harness and statistics}
A native three-process harness (orchestrator; NeuRAD renderer serving six photorealistic
cameras at the \emph{driven} pose; model node reconstructing the UniAD perception input from
rendered frames, with an in-place re-quantization endpoint) runs both benchmarks:
$16$ adversarial scene/scenario instances (frontal/side/stationary) and $14$ free-drive
scenes, $10$ seeds each per configuration ($\approx2{,}400$ rollouts across sweeps).
All aggregate numbers carry $95\%$ percentile-bootstrap confidence intervals over
scenario-level means ($10^4$ resamples); recipe comparisons use paired per-scenario
differences.

\subsection{On-device methodology (capture--replay)}
\label{sec:replay}
The fused planner input is quantization-independent (perception and embeddings are
unquantized), so we capture it on the A100, replay it through the decoder-only planner on a
Jetson AGX Orin 64\,GB at each precision, and score with the same evaluator---isolating
device$\times$quantization without porting UniAD's custom CUDA ops.

\section{Results}
\subsection{Level 1: open-loop L2 is blind}
\begin{table}[t]\centering
\caption{Open-loop nuScenes-mini ($n{=}162$), ST-P3 L2 (m).}
\label{tab:openloop}
\begin{tabular}{lcc}
\toprule
Setting & L2 (m) & well-formed \\
\midrule
FP16 & 0.33 & 100\% \\
INT8 (per-channel) & 0.33 & 100\% \\
INT4 (per-channel) & 0.33 & 100\% \\
INT3 & 1.91 & 92\% \\
INT2 & 6.45 & 0\% \\
\midrule
FP16, ego-state+history zeroed & \textbf{6.38} ($19\times$) & 100\% \\
\bottomrule
\end{tabular}
\end{table}
Weight quantization is lossless to INT4 (Table~\ref{tab:openloop}); INT3/INT2 degradation is
generation collapse, not graded planning. The ego-ablation inflates L2 $19\times$ with
perception untouched: the metric is $\approx$ ego-extrapolation and cannot see quantization
damage. Every 4-bit configuration in this paper passes this level.

\subsection{Level 2: coherence certifies group-wise INT4---incorrectly}
\label{sec:coherence}
In adversarial closed-loop rollouts, per-channel INT4 collapses generation ($\le15\%$
well-formed on 15/16 instances, mean $3.7\%$) while group-wise INT4 restores well-formedness
(mean $73.5\%$ vs FP16 $76.0\%$)---the granularity mechanism of Sec.~\ref{sec:configs}
working as predicted, and a quantization$\times$domain-shift interaction (the same
per-channel INT4 is lossless in-distribution, Table~\ref{tab:openloop}). By this metric
group-wise INT4 appears to \emph{retain closed-loop behavior}. The next two levels show that
conclusion is wrong: well-formed trajectory text is not safe driving.

\subsection{Level 3: route competence certifies AWQ---incorrectly}
\label{sec:competence}
\begin{table}[t]\centering
\caption{Free-drive route competence (clean 7 scenes $\times$ 10 seeds; artifact-corrected;
$95\%$ bootstrap CI over scenes). Success = in-corridor (cross-track $\le4$\,m) $\wedge$
moved ($\ge5$\,m) $\wedge$ no genuine collision.}
\label{tab:freedrive}
\begin{tabular}{lcc}
\toprule
Config & success & 95\% CI \\
\midrule
FP16              & 94.3\% & [82.9, 100] \\
\textbf{W8}       & \textbf{97.1\%} & [91.4, 100] \\
AWQ W4 g128       & 87.1\% & [74.3, 95.7] \\
group W4 g128     & 57.1\% & [41.4, 74.3] \\
W4A4+DCT          & 38.6\% & [10.0, 71.4] \\
naive W4          & 18.6\% & [0.0, 47.1] \\
\bottomrule
\end{tabular}
\end{table}
The free-drive benchmark (Table~\ref{tab:freedrive}) removes the freeze confound: naive W4,
whose malformed outputs freeze the ego, drives only $3.1$\,m on average and fails the
progress gate ($18.6\%$)---$45\%$ of its collisions are artifacts of being struck while
frozen, versus $5\%$ genuine. Group-wise W4 moves and tracks the route but succeeds on only
$57.1\%$ (off-route in $21\%$ of rollouts): coherent text, mediocre driving. AWQ-W4, however,
drives at near-FP16 competence ($87.1\%$, CIs overlapping FP16's). By this metric AWQ-W4
appears deployable. The final level rejects it.

\subsection{Level 4: graded safety rejects every 4-bit recipe; INT8 is certified}
\label{sec:safety}
\begin{table}[t]\centering
\caption{Graded closed-loop safety: NeuroNCAP score (0--5, impact-speed graded), 16
scene/scenario instances $\times$ 10 seeds, $95\%$ bootstrap CI over instances.}
\label{tab:safety}
\begin{tabular}{lccc}
\toprule
Config & NCAP & 95\% CI & coll.\ rate \\
\midrule
FP16              & 4.43 & [3.65, 4.93] & 12.1\% \\
\textbf{W8}       & \textbf{4.50} & [3.82, 4.94] & \textbf{10.6\%} \\
naive W4          & 2.62 & [1.60, 3.63] & 60.0\% \\
group W4 g128     & 2.62 & [1.53, 3.70] & 53.6\% \\
AWQ W4 g128       & 2.62 & [1.49, 3.71] & 51.9\% \\
W4A4+DCT          & 2.86 & [1.87, 3.82] & 47.5\% \\
\bottomrule
\end{tabular}
\end{table}
On the graded score (Table~\ref{tab:safety}), \textbf{every} 4-bit recipe collapses to
$2.6$--$2.9$---naive, group-wise, AWQ, and rotated alike---losing $\approx40\%$ of FP16's
safety score and quadrupling the collision rate. The paired group$-$naive difference is
$+0.01$ with $95\%$ CI $[-0.35,+0.34]$: granularity confers \emph{no} safety benefit despite
fully restoring coherence. AWQ's near-FP16 route competence coexists with naive-level
safety: a policy that drives well and cannot handle emergencies.
\textbf{Per-channel INT8 is statistically indistinguishable from FP16}: paired
$\Delta=-0.00$, $95\%$ CI $[-0.17,+0.20]$, with matching route competence
(Table~\ref{tab:freedrive}). W8 is the recipe this study certifies.

The 4-bit failures are \emph{structured}: all four recipes fail predominantly on
\emph{stationary}-obstacle scenarios (long-horizon stopping) while passing most frontal
ones. The mechanism is open (Sec.~\ref{sec:limits}); the structure suggests a
decision-horizon sensitivity rather than a generic accuracy loss.

\subsection{The hierarchy, assembled}
\begin{table}[t]\centering
\caption{Each metric certifies configurations the next rejects. \checkmark\ = passes at
$\ge90\%$ of FP16 (within noise for level 1).}
\label{tab:hierarchy}
\begin{tabular}{lcccc}
\toprule
Config & open-loop & coherence & competence & safety \\
\midrule
naive W4   & \checkmark & --- & --- & --- \\
group W4   & \checkmark & \checkmark & --- & --- \\
AWQ W4     & \checkmark & \checkmark & \checkmark & --- \\
\textbf{W8} & \checkmark & \checkmark & \checkmark & \checkmark \\
\bottomrule
\end{tabular}
\end{table}
Table~\ref{tab:hierarchy} is the paper in one view: open-loop certifies everything;
coherence catches naive W4 but certifies group-wise; competence catches group-wise but
certifies AWQ; only graded safety separates AWQ from W8. No cheaper prefix of the hierarchy
reaches the correct verdict.

\paragraph{Single-scene pilots also mislead.} Before full-scale evaluation, single-scene
closed-loop pilots scored group-W4, AWQ-W4, \emph{and} W4A4+DCT at a perfect $5.0$ NCAP (on
a frontal scenario each recipe ultimately passes); all three collapsed at 14-scene scale.
Closed-loop evaluation does not certify a recipe unless it also spans scenario diversity.

\subsection{On-device (Jetson AGX Orin 64\,GB)}
\label{sec:orin}
\begin{table}[t]\centering
\caption{Measured on AGX Orin (1024-token context $\to$ 60 tokens, median of 30). Latency is
precision-flat under fake-quant; the deployable win is the packed footprint.}
\label{tab:orinsys}
\begin{tabular}{lccc}
\toprule
Config & median lat. & runtime peak & packed weights \\
\midrule
FP16       & $4.02$\,s & $2.88$\,GB & $0.716$\,GB \\
\textbf{W8 g128} & $4.01$\,s & $2.88$\,GB & \textbf{$0.363$\,GB} \\
INT4 g128  & $4.02$\,s & $2.88$\,GB & $0.185$\,GB \\
\bottomrule
\end{tabular}
\end{table}
Replaying real inputs (Sec.~\ref{sec:replay}) on the Orin over all $162$ mini-val samples
reproduces the published FP16 accuracy ($0.345$\,m avg L2, vs $0.35$ reported) and matches
A100 trajectories to $\approx1.1$\,cm (FP16). Notably, the blindness of level 1 reproduces
on-device: per-sample W4 plans differ from FP16 by $14$--$24$\,cm while aggregate L2 moves
$<3$\,cm---the metric averages the damage away on the deployment hardware itself. The
certified W8 recipe packs the backbone to $0.36$\,GB ($2\times$); the whole planner node fits
in $2.88$\,GB. A packed-INT8 \emph{speedup} requires a custom kernel and is future work; the
edge bottleneck remains UniAD perception throughput, not the quantized planner.

\section{Discussion: a certification protocol}
(1) \textbf{Never certify on open-loop metrics}---they are ego-extrapolation and pass every
configuration that later fails. (2) \textbf{Do not stop at intermediate closed-loop
proxies}: well-formedness and route competence each certified a recipe that fails graded
safety. (3) \textbf{Use a graded safety score over diverse scenarios}: binary collision
saturates; single scenes certified three failing recipes. (4) \textbf{Gate competence
metrics against degenerate behavior}: without an anti-freeze gate and collision-artifact
taxonomy, a frozen policy scores as safe. (5) \textbf{Ship INT8; treat sub-8-bit as
uncertified} for driving until a recipe passes the full hierarchy---granularity, AWQ, and
rotation each fix the failure mode one level below the one that matters.

\section{Threats to Validity}
\label{sec:limits}
\textbf{Renderer domain gap:} rendered inputs are OOD to real nuScenes; claims are relative
across configurations under identical rendering. \textbf{Artifact correction:} free-drive
excludes spawn/rear collisions by a geometric classifier; 2 scenes remain unusable (ego
spawns in contact) and 1 is ambiguous, so competence CIs are computed on the clean 7 (all-14
numbers, which are consistent, are in the repository). \textbf{Fake-quant:} accuracy effects
are real; packed speedups are not yet realized (footprints are measured). \textbf{One
checkpoint, one simulator:} a 0.5B model on NeuroNCAP; the stationary-scenario failure of
4-bit recipes is characterized but not yet explained---diagnosing it, and testing whether a
mechanism-targeted mixed-precision recipe can pass the hierarchy below 8 bits, is ongoing
work. \textbf{AWQ calibration:} one rendered-domain rollout; stronger calibration could
shift AWQ's competence, though its safety failure matches recipes that need no calibration.

\section{Reproducibility}
The release includes the open-loop sweeps with ego-ablation, the three-process harness with
in-place re-quantization, the adversarial and free-drive scenario generators and runners
(\texttt{run\_freedrive.sh}, \texttt{run\_positive\_hunt.sh}), the metric tooling with
bootstrap CIs and the collision-artifact classifier
(\texttt{aggregate\_ncap\_success.py}, \texttt{eval\_freedrive\_route.py}), the on-device
suite, and all result JSONs. Every number in this paper is produced by those scripts.

\section{Conclusion}
Certifying a quantized driving VLA is a hierarchy problem: open-loop L2, generation
coherence, and route competence each certify recipes that graded closed-loop safety
rejects---group-wise INT4 restores syntax, AWQ-INT4 restores competence, and neither
restores safety, which every 4-bit recipe fails ($2.6$--$2.9$ vs FP16's $4.4$) with failures
concentrated on long-horizon stationary scenarios. Per-channel INT8 passes every level of
the hierarchy, indistinguishably from FP16, at $2\times$ compression, and reproduces on a
Jetson AGX Orin. Until a sub-8-bit recipe passes the full hierarchy, INT8 is the certified
deployment point for driving VLAs---and no metric short of graded, scenario-diverse,
closed-loop safety testing should be trusted to say otherwise.

\begin{thebibliography}{9}
\bibitem{opendrivevla} X.\ Zhou et al. \emph{OpenDriveVLA: Towards End-to-end Autonomous
Driving with Large Vision Language Action Model}. arXiv:2503.23463, 2025.
\bibitem{uniad} Y.\ Hu et al. \emph{Planning-oriented Autonomous Driving (UniAD)}. CVPR 2023.
\bibitem{vad} B.\ Jiang et al. \emph{VAD: Vectorized Scene Representation for Efficient
Autonomous Driving}. ICCV 2023.
\bibitem{egostatus} Z.\ Li et al. \emph{Is Ego Status All You Need for Open-Loop End-to-End
Autonomous Driving?} CVPR 2024. arXiv:2312.03031.
\bibitem{neuroncap} W.\ Ljungbergh et al. \emph{NeuroNCAP: Photorealistic Closed-Loop Safety
Testing for Autonomous Driving}. ECCV 2024. arXiv:2404.07762.
\bibitem{neurad} A.\ Tonderski et al. \emph{NeuRAD: Neural Rendering for Autonomous Driving}.
CVPR 2024.
\bibitem{awq} J.\ Lin et al. \emph{AWQ: Activation-aware Weight Quantization for On-Device
LLM Compression and Acceleration}. MLSys 2024. arXiv:2306.00978.
\bibitem{gptq} E.\ Frantar et al. \emph{GPTQ: Accurate Post-Training Quantization for
Generative Pre-trained Transformers}. ICLR 2023. arXiv:2210.17323.
\bibitem{llmint8} T.\ Dettmers et al. \emph{LLM.int8(): 8-bit Matrix Multiplication for
Transformers at Scale}. NeurIPS 2022. arXiv:2208.07339.
\bibitem{smoothquant} G.\ Xiao et al. \emph{SmoothQuant: Accurate and Efficient Post-Training
Quantization for Large Language Models}. ICML 2023. arXiv:2211.10438.
\bibitem{quarot} S.\ Ashkboos et al. \emph{QuaRot: Outlier-Free 4-Bit Inference in Rotated
LLMs}. NeurIPS 2024. arXiv:2404.00456.
\end{thebibliography}
\end{document}
